\section{The binomial distribution}

We say that $X$ has the binomial distribution with
parameters $n\ge 1$ and $0<p<1$, and denote
$X\sim \operatorname{Bin}(n,p)$, if $X$ satisfies

\[
\mathrm{P}(X = k) = \binom{n}{k} p^k (1-p)^{n-k} \qquad (0\le k\le n).
\]

\noindent $\textbf{Properties:}$

\begin{enumerate}[label=\arabic*.]
\item $\operatorname{Bin}(1,p) \stackrel{d}{=}
\operatorname{Ber}(p)$.

\item $\operatorname{Bin}(n,p) \boxplus
\operatorname{Bin}(m,p) \stackrel{d}{=}
\operatorname{Bin}(n+m,p)$.

\item $\operatorname{Bin}(n,p) \stackrel{d}{=}
\mathop{\boxplus}\limits_{k=1}^n
\operatorname{Ber}(p)$.
That is, the binomial distribution models the number
of successes when $n$ identical experiments are
performed independently, where each experiment has
probability $p$ of success.

\item If $X\sim \operatorname{Bin}(n,p)$ then
$\mathbf{E} X = np$, $\mathbf{V}(X) = np(1-p)$.
\end{enumerate}
